% 后处理
\makeatletter
\define@key{tongji}{cheadingtitlenew}{\cheadingtitlenew{#1}}
\gdef\cheadingtitlenew#1{\gdef\tongji@cheadingtitlenew{#1}}
\tongjideletelinebreakcmd{tongji@ctitle}{tongji@posttitle}
\cheadingtitlenew{\tongji@posttitle}

% === 之前修改的英文摘要页眉补丁 ===
\patchcmd{\tongji@makeabstract}
  {\xparameter@heading@eother{\tongji@eabstractname}}
  {\xparameter@heading@eother{Abstract}}
  {}{}

% ===============================================
% === 新增：生成单独书脊页的命令 ===
\newcommand{\makespine}{
  \clearpage
  \thispagestyle{empty}
  % 页面总高 29.7cm，减去上下各 5cm，中间文本区高度设为 19.7cm。
  % 利用 \vspace*{\fill} 上下弹性挤压，完美实现物理边距各 5cm。
  \vspace*{\fill}
  \begin{center}
    % 宽度设为 1.2em，强制每个汉字自动换行，实现垂直排版
    \parbox[c][19.7cm][s]{1.2em}{
      \centering
      \fangsong\zihao{4}\bfseries         % 仿宋、四号、加粗
      \fontsize{14bp}{16bp}\selectfont    % 四号字对应14bp，行距严格指定为16bp(磅)
      \setlength{\parskip}{0pt}           % 段前段后 0 磅
      
      \tongji@posttitle                   % 使用无换行符版本的标题
      退化条件下的鲁棒语音转换算法研究
      \par\vfill
      \tongji@cauthor  
      宋昕宁                              % 姓名
      \par\vfill
      同济大学
    }
  \end{center}
  \vspace*{\fill}
  \clearpage
}
% ===============================================
\makeatother